(1) 圆形瞳函数的二维Fourier-Bessel变换给远场振幅 $A(u)\propto2J_1(u)/u$，其中 $u=\pi D\theta/\lambda$ 且 $u\to0$ 时比值趋一。于是相对中心峰强度 $I(\theta)/I(0)=[2J_1(u)/u]^2$，这是经典Airy图样。
(2) 第一暗环由 $J_1(u_1)=0$ 决定，其首个正根 $u_1\simeq3.831706$，所以 $\theta_1=u_1\lambda/(\pi D)\simeq1.22\,\lambda/D$。它是角半径而不是直径，适用于足够小衍射角下的近轴近似。
(3) 归一包围能量分数可写 $F(u)=\frac12\int_0^u [2J_1(s)/s]^2s\,ds$。用Bessel恒等式 $d[J_0^2+J_1^2]/du=-2J_1(u)^2/u$，得 $F(u)=1-J_0(u)^2-J_1(u)^2$。在第一暗环 $J_1(u_1)=0$，故 $F(u_1)=1-J_0(u_1)^2\simeq0.838$，即约83.8%。
(4) 角半径与 $\lambda/D$ 成正比，因此大孔径、小波长使中央斑更窄。圆形瞳孔因旋转对称而含Bessel函数，矩形瞳孔的二维Fourier变换则为不同方向正弦函数比值的乘积；不能混用第一暗纹的常数系数。
